\documentclass[11pt]{amsproc}

\usepackage[a4paper,margin=2.6cm]{geometry}
\usepackage{amsmath,amssymb,amsthm}
\usepackage[T1]{fontenc}
\usepackage[utf8]{inputenc}

\newcommand{\id}{\mathrm{id}}
\newcommand{\mcm}{\operatorname{mcm}}

\begin{document}

\section*{Problema}

En una pizarra hay una tabla con dos filas y $n\ge 3$ columnas, numeradas
del $1$ al $n$. En cada fila escribimos los números del $1$ al $n$, en algún
orden.

La fila de arriba no cambia nunca. En cada paso, debajo de cada número de la
fila de arriba escribimos el número de la columna en la que ese número estaba
en la fila de abajo antes del paso. Por ejemplo, si $n=3$, la fila de arriba
es $2\,3\,1$ y la de abajo es $1\,3\,2$, después de un paso la fila de abajo
es $3\,2\,1$: el $2$ estaba en la columna $3$, el $3$ en la columna $2$ y el
$1$ en la columna~$1$.

{\itshape Hallar el menor entero positivo $N$ tal que, para cualquier
elección de las dos filas iniciales, después de $N$ pasos la fila de abajo
vuelve a ser la del principio.}

\section*{Solución}

Pensemos cada fila como una permutación de $\{1,2,\dots,n\}$: $s(j)$ es el
número que está en la columna $j$ de la fila de arriba, y $f(j)$ el de la
fila de abajo. Escribimos $gh=g\circ h$, $g^{-1}$ para la inversa de $g$ e
$\id$ para la identidad. El número $s(j)$ está en la columna $f^{-1}(s(j))$
de la fila de abajo. Entonces un paso cambia $f$ por $f^{-1}s$.

Vamos a ver que la respuesta es $N=2L$, donde $L=\mcm(1,2,\dots,n)$.

\subsection*{Paso 1: dos pasos}

Después de dos pasos, la fila de abajo es
\[
  (f^{-1}s)^{-1}s=s^{-1}fs .
\]
O sea, dos pasos conjugan $f$ por $s$. Como $s$ no cambia, después de $2k$
pasos la fila de abajo es $s^{-k}fs^k$, y después de $2k+1$ pasos es
\[
  \bigl(s^{-k}fs^k\bigr)^{-1}s=s^{-k}f^{-1}s^{k+1}.
\]

\subsection*{Paso 2: las cantidades impares no sirven}

Supongamos que las dos filas empiezan iguales, o sea, $f=s$. Después de
$2k+1$ pasos, la fila de abajo es $s^{-k}s^{-1}s^{k+1}=\id$, es decir,
$1\,2\,\dots\,n$. Para el primer paso esto se ve sin cuentas: si las dos
filas son iguales, cada número queda en su propia columna. Si elegimos las dos filas iguales y
distintas de $1\,2\,\dots\,n$, la fila de abajo no vuelve a la inicial
después de una cantidad impar de pasos.

\subsection*{Paso 3: las cantidades pares}

Por el paso 1, después de $2k$ pasos la fila de abajo vuelve a ser la inicial
para cualquier elección de las filas si y solo si $s^{-k}fs^k=f$, es decir,
$s^kf=fs^k$, para todas las permutaciones $s$ y $f$. Si $s^k=\id$ para toda
$s$, esto vale. Al revés, si vale, entonces cada $s^k$ conmuta con todas las
permutaciones.

Pero la única permutación que conmuta con todas es la identidad, y acá se usa
que $n\ge 3$. Sea $\pi$ una permutación que conmuta con todas, y sea $i$ un
elemento. Para cada $j\ne i$, sea $t$ la permutación que intercambia $i$ y
$j$ y deja fijo todo lo demás. Como $\pi t=t\pi$, tenemos $\pi t\pi^{-1}=t$.
Pero $\pi t\pi^{-1}$ intercambia $\pi(i)$ y $\pi(j)$ y deja fijo todo lo
demás. Por eso $\{\pi(i),\pi(j)\}=\{i,j\}$, y en particular
$\pi(i)\in\{i,j\}$. Como $n\ge 3$, hay dos elementos distintos $j$ y $j'$,
ambos distintos de $i$, así que $\pi(i)\in\{i,j\}\cap\{i,j'\}=\{i\}$.
Entonces $\pi=\id$.

Así que $2k$ pasos sirven si y solo si $s^k=\id$ para toda permutación $s$.
Si $s$ tiene un ciclo de longitud $\ell$, en ese ciclo $s^k$ avanza $k$
lugares, y eso es la identidad justo cuando $\ell$ divide a $k$. Para cada
$\ell\le n$ hay permutaciones con un ciclo de longitud $\ell$. Por eso
$s^k=\id$ para toda $s$ si y solo si todos los números $1,2,\dots,n$ dividen
a $k$, es decir, si y solo si $L$ divide a $k$.

En resumen, los $N$ que sirven son los de la forma $2k$ con $k$ múltiplo de
$L$, y el menor es $N=2L=2\mcm(1,2,\dots,n)$.

\section*{Observaciones}

\begin{enumerate}
  \item Algunos valores: para $n=3,4,5,6,7$ la respuesta es $12$, $24$,
        $120$, $120$ y $840$. Lo verificamos por computadora para
        $3\le n\le 6$, simulando la pizarra con todas las filas iniciales
        posibles.

  \item Para $n=2$ la respuesta no es $2\mcm(1,2)=4$ sino $2$, porque ahí
        todas las permutaciones conmutan y dos pasos siempre devuelven la
        fila inicial. Por eso el enunciado pide $n\ge 3$.

  \item Para una versión con pista, se puede pedir primero describir la fila
        de abajo después de dos pasos. Ese es el paso 1, y con eso el resto
        es bastante directo.

  \item Es una versión del problema de la función
        $F(\sigma,\tau)=(\tau,\tau^{-1}\sigma\tau)$ en pares de
        permutaciones. Esa $F$ no cambia el producto $c=\sigma\tau$, y la
        primera coordenada pasa de $\sigma$ a $\sigma^{-1}c$. Acá la fila de
        arriba hace de $c$ y la de abajo hace de $\sigma$, así que $F^N$ es
        la identidad si y solo si $N$ pasos devuelven siempre la fila de
        abajo inicial. La pizarra se queda solamente con lo que importa: el
        producto fijo y una coordenada.

  \item El problema sale del teorema~3.13 de A.~Smoktunowicz y
        L.~Vendramin, \emph{On skew braces}, arXiv:1705.06958. El teorema
        dice que si $A$ es un skew brace finito con más de un elemento,
        entonces la solución $r_A$ de la ecuación de Yang--Baxter asociada a
        $A$ tiene orden $2d$, donde $d$ es el exponente del grupo aditivo
        módulo su centro. La función $F$ de arriba es $r_A$ para el skew
        brace trivial sobre el grupo simétrico, y para $n\ge 3$ el centro es
        trivial, así que $d=\mcm(1,2,\dots,n)$. El paso 2 es el mismo
        argumento del paper para descartar las potencias impares, y el
        paso 1 corresponde a su lema~3.12.
\end{enumerate}

\end{document}
